% kacho E2E Tier-4 fixture: academic-style article with sections, floats
% (\rule placeholders, not images), tables, math and footnotes.
\documentclass[11pt]{article}
\usepackage{geometry}
\usepackage{amsmath,amssymb}
\usepackage{lipsum}
\usepackage[theme=haru]{kacho}
\title{Procedural Brushwork for Flower-and-Bird Painting}
\author{A.~Author \and B.~Author}
\date{1 November 2026}
\begin{document}
\maketitle
\begin{abstract}
We describe a procedural model of tapered brush strokes and layered washes.
\lipsum[1][1-4]
\end{abstract}
\section{Introduction}
Kach\=o-e prints\footnote{Literally ``flower-and-bird pictures''.} combine
botanical motifs with birds. \lipsum[2]
\section{Model}
A stroke of length $L$ has width $w(t)=w_0\,\sin(\pi t)^{0.6}$ for
$t\in[0,1]$, and its centre line is a cubic B\'ezier
\begin{equation}
  B(t)=(1-t)^3P_0+3(1-t)^2tP_1+3(1-t)t^2P_2+t^3P_3 .
\end{equation}
Layered washes use $n\ge 6$ layers of opacity $\alpha\le 0.12$, giving a
total coverage of
\begin{align}
  A_n &= 1-(1-\alpha)^n, & \lim_{n\to\infty} A_n &= 1.
\end{align}
\lipsum[3]
\begin{figure}[t]
  \centering
  \rule{0.8\linewidth}{4cm}
  \caption{A placeholder figure drawn with a rule.}\label{fig:a}
\end{figure}
\begin{table}[h]
  \centering
  \begin{tabular}{lrr}\hline
  Motif & Strokes & Layers\\\hline
  Ume & 120 & 6\\ Bamboo & 80 & 8\\ Iris & 64 & 7\\\hline
  \end{tabular}
  \caption{Stroke counts.}
\end{table}
\section{Results}
See Figure~\ref{fig:a}.\footnote{A second footnote with \emph{emphasis}.}
\lipsum[4-7]
\begin{figure}[p]
  \centering
  \rule{0.9\linewidth}{0.5\textheight}
  \caption{A full-page float placeholder.}
\end{figure}
\section{Conclusion}
\lipsum[8-10]
\end{document}
